\section{Background}
\label{sec:background}

\subsection{State-Space Models and Bayesian Filtering}
\label{subsec:state-space-models}

A discrete-time state-space model consists of an unobserved state
\(X_t\) and an observation \(Y_t\),

\begin{align}
	X_t &\sim p(x_t\mid x_{t-1}), \\
	Y_t &\sim p(y_t\mid x_t).
\end{align}

The filtering problem is to compute the conditional distribution

\begin{equation}
	p(x_t\mid y_{0:t}),
\end{equation}

where \(y_{0:t}\) denotes all observations available through time \(t\).

The recursion consists of a prediction step,

\begin{equation}
	p(x_t\mid y_{0:t-1})
	=
	\int
	p(x_t\mid x_{t-1})
	p(x_{t-1}\mid y_{0:t-1})
	\,dx_{t-1},
\end{equation}

followed by a correction after observing \(y_t\),

\begin{equation}
	p(x_t\mid y_{0:t})
	\propto
	p(y_t\mid x_t)
	p(x_t\mid y_{0:t-1}).
\end{equation}


\subsection{Kalman Filtering}
\label{subsec:kalman-filter}

For a linear-Gaussian model,

\begin{align}
	X_t &= A X_{t-1} + \varepsilon_t,
	& \varepsilon_t &\sim \mathcal{N}(0,Q), \\
	Y_t &= C X_t + \eta_t,
	& \eta_t &\sim \mathcal{N}(0,R),
\end{align}

the filtering distribution remains Gaussian. It is therefore sufficient
to update its mean and variance.

If \(m_t^{-}\) and \(P_t^{-}\) denote the predictive mean and variance,
the scalar Kalman update is

\begin{align}
	K_t &=
	\frac{P_t^- C}{C^2P_t^-+R}, \\
	m_t &=
	m_t^- + K_t(Y_t-Cm_t^-), \\
	P_t &=
	(1-K_tC)P_t^-,
\end{align}

with prediction

\begin{align}
	m_{t+1}^- &= A m_t, \\
	P_{t+1}^- &= A^2P_t + Q.
\end{align}

For the linear-Gaussian model used in Stage~I, these recursions provide
the exact filtering mean and therefore a convenient benchmark.


\subsection{Particle Filtering}
\label{subsec:particle-filtering}

A particle filter approximates the filtering distribution using weighted
samples,

\begin{equation}
	p(x_t\mid y_{0:t})
	\approx
	\sum_{i=1}^{N}
	w_t^{(i)}
	\delta_{X_t^{(i)}}(x_t),
\end{equation}

with normalized weights satisfying

\begin{equation}
	\sum_{i=1}^{N} w_t^{(i)} = 1.
\end{equation}

In a bootstrap particle filter, particles are propagated using the state
transition,

\begin{equation}
	X_t^{(i)}
	\sim
	p(x_t\mid X_{t-1}^{(i)}),
\end{equation}

and weighted according to the observation likelihood,

\begin{equation}
	\widetilde{w}_t^{(i)}
	=
	p(y_t\mid X_t^{(i)}).
\end{equation}

After normalization,

\begin{equation}
	w_t^{(i)}
	=
	\frac{\widetilde{w}_t^{(i)}}
	{\sum_{j=1}^{N}\widetilde{w}_t^{(j)}},
\end{equation}

quantities such as the filtered mean can be estimated by

\begin{equation}
	\widehat{X}_t
	=
	\sum_{i=1}^{N} w_t^{(i)}X_t^{(i)}.
\end{equation}

The particle population is then resampled according to the normalized
weights. Since the method is Monte Carlo, its statistical error is
expected to decrease at the usual \(O(N^{-1/2})\) rate under standard
conditions.


\subsection{Parallel Structure}
\label{subsec:pf-parallel-structure}

Particle propagation and likelihood evaluation are independent across
particles and are therefore natural targets for data parallelism.
Normalization and posterior estimation require reductions over the
particle population, while resampling depends on the full set of
weights.

The distinction between particle-local work and global operations is the
main computational issue considered in Stage~I.